% GENERATED FILE. Edit canonical lessons, then run npm run generate:books.
% canonical-source-hash: fnv1a64:058cb24cf643f288
% canonical-lessons: TE-C478-alpapidanam, TE-C478-penugali, TE-C478-sudigali, TE-C478-galivana, TE-C478-balahinapadu

\chapter{Storms}
\label{ch:te-storms}
\hlchaptermodality{478}

\begin{chapteropening}
\textbf{By the end of this chapter:} \emph{I can say a low-pressure area, a gale, a whirlwind, a rainstorm and to lose strength.}

Complete the last lesson of chapter 478: I can say a low-pressure area, a gale, a whirlwind, a rainstorm and to lose strength.
\end{chapteropening}

\section[\texorpdfstring{\te{అల్పపీడనం} (alpapīḍanaṁ)}{అల్పపీడనం (alpapīḍanaṁ)}]{\te{అల్పపీడనం} (alpapīḍanaṁ) — a low-pressure area}

\label{lesson:TE-C478-alpapidanam}



\begin{quote}
Before the new one: say the Telugu for a shower of rain, then the Telugu for to get soaked.
\end{quote}

\subsection*{You'll want to know: \te{అల్పపీడనం}}
\textbf{\te{అల్పపీడనం}} — \emph{alpapīḍanaṁ} — \textquotedblleft{}a low-pressure area\textquotedblright{}.

\begin{quote}
\te{సముద్రంలో} \te{అల్పపీడనం} \te{ఏర్పడింది}.
\end{quote}

\begin{quote}
\emph{samudraṁlō alpapīḍanaṁ ērpaḍindi.}
\end{quote}

\begin{quote}
A low-pressure area has formed over the sea.
\end{quote}

The weather news says it every monsoon: \te{అల్ప}, low, + \te{పీడనం}, pressure. A storm often follows.

\begin{cousinweb}
\textbf{\te{అల్పపీడనం}} (\emph{alpapīḍanaṁ}) is a modern weather word built from two Sanskrit words: \emph{alpa}, \textquotedblleft{}small, little\textquotedblright{}, and \emph{pīḍana}, \textquotedblleft{}pressing, squeezing\textquotedblright{} (Monier-Williams).
\end{cousinweb}

\begin{grammarlens}[title={What you've built}]
One more word for getting around.
\end{grammarlens}

\subsection*{Your turn}
\begin{itemize}
\raggedright
  \item \emph{Say it:} \emph{alpapīḍanaṁ}
  \item \emph{Say it:} \emph{alpapīḍanaṁ}, once more
  \item \emph{Say it:} \emph{samudraṁlō alpapīḍanaṁ ērpaḍindi.} — A low-pressure area has formed over the sea.
  \item \emph{Recall:} say the Telugu for a shower of rain, then the Telugu for to get soaked, then say \emph{alpapīḍanaṁ} again
\end{itemize}

\begin{tcolorbox}[breakable,colback=teal!4,colframe=teal!35!black,title={Before you move on}]
What does \te{అల్పపీడనం} mean? (\textquotedblleft{}A low-pressure area\textquotedblright{}.) Say it once more.
\end{tcolorbox}
\section[\texorpdfstring{\te{పెనుగాలి} (penugāli)}{పెనుగాలి (penugāli)}]{\te{పెనుగాలి} (penugāli) — a gale; a strong wind}

\label{lesson:TE-C478-penugali}



\begin{quote}
Before the new one: say the Telugu for to get soaked, then the Telugu for a low-pressure area.
\end{quote}

\subsection*{You'll want to know: \te{పెనుగాలి}}
\textbf{\te{పెనుగాలి}} — \emph{penugāli} — \textquotedblleft{}a gale; a strong wind\textquotedblright{}.

\begin{quote}
\te{పెనుగాలికి} \te{చాలా} \te{చెట్లు} \te{కూలాయి}.
\end{quote}

\begin{quote}
\emph{penugāliki cālā ceṭlu kūlāyi.}
\end{quote}

\begin{quote}
Many trees came down in the gale.
\end{quote}

\te{పెను}, great, + \te{గాలి}, wind.

\begin{cousinweb}
\textbf{\te{పెనుగాలి}} (\emph{penugāli}) is native Dravidian: \te{పెను}, \textquotedblleft{}great, big\textquotedblright{}, which DEDR 4411 lists beside Tamil \emph{peru}, \textquotedblleft{}great\textquotedblright{}, and \te{గాలి}, \textquotedblleft{}wind\textquotedblright{} (DEDR 1499).
\end{cousinweb}

\begin{grammarlens}[title={What you've built}]
One more word for getting around.
\end{grammarlens}

\subsection*{Your turn}
\begin{itemize}
\raggedright
  \item \emph{Say it:} \emph{penugāli}
  \item \emph{Say it:} \emph{penugāli}, once more
  \item \emph{Say it:} \emph{penugāliki cālā ceṭlu kūlāyi.} — Many trees came down in the gale.
  \item \emph{Recall:} say the Telugu for to get soaked, then the Telugu for a low-pressure area, then say \emph{penugāli} again
\end{itemize}

\begin{tcolorbox}[breakable,colback=teal!4,colframe=teal!35!black,title={Before you move on}]
What does \te{పెనుగాలి} mean? (\textquotedblleft{}A gale; a strong wind\textquotedblright{}.) Say it once more.
\end{tcolorbox}
\section[\texorpdfstring{\te{సుడిగాలి} (suḍigāli)}{సుడిగాలి (suḍigāli)}]{\te{సుడిగాలి} (suḍigāli) — a whirlwind}

\label{lesson:TE-C478-sudigali}



\begin{quote}
Before the new one: say the Telugu for a low-pressure area, then the Telugu for a gale.
\end{quote}

\subsection*{You'll want to know: \te{సుడిగాలి}}
\textbf{\te{సుడిగాలి}} — \emph{suḍigāli} — \textquotedblleft{}a whirlwind\textquotedblright{}.

\begin{quote}
\te{సుడిగాలికి} \te{ఇంటి} \te{పైకప్పు} \te{ఎగిరిపోయింది}.
\end{quote}

\begin{quote}
\emph{suḍigāliki iṇṭi paikappu egiripōyindi.}
\end{quote}

\begin{quote}
In the whirlwind the roof of the house blew off.
\end{quote}

\te{సుడి}, a whirl, + \te{గాలి}, wind. A whirlpool in water is a \te{సుడిగుండం}.

\begin{cousinweb}
\textbf{\te{సుడిగాలి}} (\emph{suḍigāli}) is native Dravidian: DEDR 2698 lists Telugu \emph{suḍi}, \textquotedblleft{}a whirl, a whirlpool, a whirlwind\textquotedblright{}, here joined to \te{గాలి}, \textquotedblleft{}wind\textquotedblright{} (DEDR 1499).
\end{cousinweb}

\begin{grammarlens}[title={What you've built}]
One more word for getting around.
\end{grammarlens}

\subsection*{Your turn}
\begin{itemize}
\raggedright
  \item \emph{Say it:} \emph{suḍigāli}
  \item \emph{Say it:} \emph{suḍigāli}, once more
  \item \emph{Say it:} \emph{suḍigāliki iṇṭi paikappu egiripōyindi.} — In the whirlwind the roof of the house blew off.
  \item \emph{Recall:} say the Telugu for a low-pressure area, then the Telugu for a gale, then say \emph{suḍigāli} again
\end{itemize}

\begin{tcolorbox}[breakable,colback=teal!4,colframe=teal!35!black,title={Before you move on}]
What does \te{సుడిగాలి} mean? (\textquotedblleft{}A whirlwind\textquotedblright{}.) Say it once more.
\end{tcolorbox}
\section[\texorpdfstring{\te{గాలివాన} (gālivāna)}{గాలివాన (gālivāna)}]{\te{గాలివాన} (gālivāna) — a rainstorm}

\label{lesson:TE-C478-galivana}



\begin{quote}
Before the new one: say the Telugu for a gale, then the Telugu for a whirlwind.
\end{quote}

\subsection*{You'll want to know: \te{గాలివాన}}
\textbf{\te{గాలివాన}} — \emph{gālivāna} — \textquotedblleft{}a rainstorm\textquotedblright{}.

\begin{quote}
\te{నిన్న} \te{రాత్రి} \te{పెద్ద} \te{గాలివాన} \te{వచ్చింది}.
\end{quote}

\begin{quote}
\emph{ninna rātri pedda gālivāna vaccindi.}
\end{quote}

\begin{quote}
There was a big rainstorm last night.
\end{quote}

\te{గాలి}, wind, + \te{వాన}, rain: wind and rain together.

\begin{cousinweb}
\textbf{\te{గాలివాన}} (\emph{gālivāna}) is native Dravidian: \te{గాలి}, \textquotedblleft{}wind\textquotedblright{} (DEDR 1499), and \te{వాన}, \textquotedblleft{}rain\textquotedblright{} (DEDR 5381).
\end{cousinweb}

\begin{grammarlens}[title={What you've built}]
One more word for getting around.
\end{grammarlens}

\subsection*{Your turn}
\begin{itemize}
\raggedright
  \item \emph{Say it:} \emph{gālivāna}
  \item \emph{Say it:} \emph{gālivāna}, once more
  \item \emph{Say it:} \emph{ninna rātri pedda gālivāna vaccindi.} — There was a big rainstorm last night.
  \item \emph{Recall:} say the Telugu for a gale, then the Telugu for a whirlwind, then say \emph{gālivāna} again
\end{itemize}

\begin{tcolorbox}[breakable,colback=teal!4,colframe=teal!35!black,title={Before you move on}]
What does \te{గాలివాన} mean? (\textquotedblleft{}A rainstorm\textquotedblright{}.) Say it once more.
\end{tcolorbox}
\section[\texorpdfstring{\te{బలహీనపడు} (balahīnapaḍu)}{బలహీనపడు (balahīnapaḍu)}]{\te{బలహీనపడు} (balahīnapaḍu) — to lose strength; to weaken}

\label{lesson:TE-C478-balahinapadu}



\begin{quote}
Before the new one: say the Telugu for a whirlwind, then the Telugu for a rainstorm.
\end{quote}

\subsection*{You'll want to know: \te{బలహీనపడు}}
\textbf{\te{బలహీనపడు}} — \emph{balahīnapaḍu} — \textquotedblleft{}to lose strength; to weaken\textquotedblright{}.

\begin{quote}
\te{తీరం} \te{దాటిన} \te{తరువాత} \te{తుపాను} \te{బలహీనపడింది}.
\end{quote}

\begin{quote}
\emph{tīraṁ dāṭina taruvāta tupānu balahīnapaḍindi.}
\end{quote}

\begin{quote}
After crossing the coast the storm weakened.
\end{quote}

Past (I …): \te{బలహీనపడ్డాను}, \emph{balahīnapaḍḍānu}.

\te{బలం}, strength, + \te{హీన}, lacking, + \te{పడు}. Said of storms, of the body and of a currency.

\begin{cousinweb}
\textbf{\te{బలహీనపడు}} (\emph{balahīnapaḍu}) is a Telugu verb built from Sanskrit \emph{bala}, \textquotedblleft{}strength\textquotedblright{}, and \emph{hīna}, \textquotedblleft{}lacking, deprived of\textquotedblright{} (Monier-Williams), with the native \te{పడు}, \textquotedblleft{}to undergo\textquotedblright{} (DEDR 3853).
\end{cousinweb}

\begin{grammarlens}[title={What you've built}]
One more word for getting around.
\end{grammarlens}

\subsection*{Your turn}
\begin{itemize}
\raggedright
  \item \emph{Say it:} \emph{balahīnapaḍu}
  \item \emph{Say it:} \emph{balahīnapaḍu}, once more
  \item \emph{Say it:} \emph{tīraṁ dāṭina taruvāta tupānu balahīnapaḍindi.} — After crossing the coast the storm weakened.
  \item \emph{Recall:} say the Telugu for a whirlwind, then the Telugu for a rainstorm, then say \emph{balahīnapaḍu} again
\end{itemize}

\begin{tcolorbox}[breakable,colback=teal!4,colframe=teal!35!black,title={Before you move on}]
What does \te{బలహీనపడు} mean? (\textquotedblleft{}To lose strength; to weaken\textquotedblright{}.) Say it once more.
\end{tcolorbox}
